\clearpage
\section{Complete recursive translation and the original proposition graphs}
\label{sec:complete-translation}
The Figure~16 graph construction is now implemented for the entire scoped
term datatype, including application and every binary example operation.
Both unchanged programs of Proposition~12.2 are translated in arbitrary
variable and atom environments. The construction retains all canonical
trees, weakenings, instances, boxes, and interface permutations prescribed
by the paper.

This completes construction of the term translation, not the observational
equivalence theorem. The translation's renaming and substitution laws beyond
its depth calculation, translation of contexts and their composition, the
copying/store machine, and the template robustness arguments remain open.
The well-formedness result for the two proposition graphs is not a replacement
for Proposition~12.2.

\subsection{The middle-block permutation}
\termdef{middle-block-permutation}
The parallel subterms expose variables and atoms separately on each side:
$(\Gamma,\Delta,\Gamma,\Delta)$. Figure~16 exchanges the middle blocks to
obtain $(\Gamma,\Gamma,\Delta,\Delta)$ before the two-use distributors.
We implement the more general typed permutation
$(A,B,C,D)\mapsto(A,C,B,D)$, preserving each local port index.

First express concatenated ports by a disjoint sum: left ports have their
original index, and right ports are offset by the left interface length.
The existing ordered-tensor join and split maps prove the inverse laws.
The two label lemmas state that looking up a left or right joined port
returns its original wire label. Empty interfaces are included.

\lean{Graph/InterfaceShuffle}{9}{21}{InterfaceShuffle.ports}

\lean{Graph/InterfaceShuffle}{23}{25}{InterfaceShuffle.left_label}

\lean{Graph/InterfaceShuffle}{27}{29}{InterfaceShuffle.right_label}

On four disjoint summands, interchange sends an $A$ index to $A$, a $B$
index to the third block, a $C$ index to the second, and a $D$ index to $D$.
Its inverse performs the reverse case distinction. Both inverse equations
are immediate by the four cases. The original-port bijection groups the
summands as $(A+B)+(C+D)$. The shuffled-port bijection first interchanges
the summands and then groups them as $(A+C)+(B+D)$. Composing the inverse
of the original bijection with the shuffled bijection gives the required
permutation of finite interface positions.

\lean{Graph/InterfaceShuffle}{32}{44}{InterfaceShuffle.interchange}

\lean{Graph/InterfaceShuffle}{47}{48}{InterfaceShuffle.original}

\lean{Graph/InterfaceShuffle}{51}{53}{InterfaceShuffle.shuffled}

\lean{Graph/InterfaceShuffle}{56}{58}{InterfaceShuffle.bijection}

For each of the four component cases, repeated application of the left
and right label lemmas shows that the original and shuffled positions have
the same wire label. To prove this for an arbitrary original interface
position, take its inverse image under the original-port bijection and
use its right inverse equation. Thus the permutation is type preserving,
not merely a bijection between interfaces of the same size.

\lean{Graph/InterfaceShuffle}{60}{65}{InterfaceShuffle.component_label}

\lean{Graph/InterfaceShuffle}{68}{71}{InterfaceShuffle.label}

\subsection{Application and binary-operation clauses}

The pair boundary leaves the two term inputs in left-to-right order.
For outputs it first concatenates the two subterm interfaces, then applies
the middle-block permutation with $A=C=\star^k$ and $B=D=\diamond^h$.
Type preservation is the composition of the tensor and permutation label
lemmas. Parallel assembly therefore constructs the actual two-subterm graph
with input $\star\otimes\star$ and grouped outputs
$\star^{2k}\otimes\diamond^{2h}$. The final cast uses only equalities of
replicate-list lengths and concatenations; it does not perform another
permutation or simplify an edge.

\lean{Copying/BinaryTranslation}{11}{22}{BinaryTranslation.pairBoundary}

\lean{Copying/BinaryTranslation}{25}{29}{BinaryTranslation.pair}

For application, serially assemble the application corolla, the paired
subterms, and $D^{\star}_{k,2}\otimes D^{\diamond}_{h,2}$. For a binary
operation use the corresponding operation corolla with the same wiring.
This covers addition, subtraction, assignment, and atom equality. These
are precisely the two right-hand constructions in Figure~16. In particular,
both full ambient environments occur in the subterm graphs before sharing;
the construction does not prune unused occurrences or weakenings.

\lean{Copying/BinaryTranslation}{32}{35}{BinaryTranslation.application}

\lean{Copying/BinaryTranslation}{38}{41}{BinaryTranslation.binary}

The pair's edge labels come from exactly one of its two subterms, so their
global language predicates imply the pair's predicate. Interface casting
preserves that predicate. Serial closure, the permitted operation label,
and the proved distributor predicate establish the language property for
both application and binary operations. These proofs inspect all nested
box contents through the existing recursive label predicate.

\lean{Copying/BinaryTranslation}{43}{49}{BinaryTranslation.pair_language}

\lean{Copying/BinaryTranslation}{51}{54}{BinaryTranslation.application_language}

\lean{Copying/BinaryTranslation}{56}{59}{BinaryTranslation.binary_language}

\subsection{Common depth allowance without graph simplification}
\termdef{translation-depth-padding}
A finite representation records a maximum box allowance. Two subterms of
an application need not have the same allowance, although they must be
assembled in one alphabet. Padding raises the smaller allowance while
preserving the graph incidence and recursively represented box contents.
It adds no term constructor, sharing edge, or box.

A depth cast changes only provably equal numeric indices, and equality
elimination proves preservation and reflection of all global label
predicates. Repeated padding is defined by iterating the earlier one-level
raising operation. Induction on the number of iterations proves label
invariance, using the previously proved raising equivalence at each step.
To pad a graph from allowance $d$ to $n\ge d$, iterate $n-d$ times and
cast along $d+(n-d)=n$. This gives padding to any chosen common allowance
with the same global label properties.

\lean{Copying/Translation}{10}{10}{Translation.castDepth}

\lean{Copying/Translation}{12}{13}{Translation.allLabels_castDepth}

\lean{Copying/Translation}{16}{18}{Translation.pad}

\lean{Copying/Translation}{20}{24}{Translation.allLabels_pad}

\lean{Copying/Translation}{27}{28}{Translation.padTo}

\lean{Copying/Translation}{30}{32}{Translation.allLabels_padTo}

One-level raising leaves the outer edge enumeration unchanged, by cases
on the original depth. Induction proves the same count equality for any
number of padding steps. A depth equality cast also preserves that count,
by equality elimination. Combining these results proves that padding to
a common allowance allocates no additional outer edge. These counts are
structural checks, not execution costs.

\lean{Copying/Translation}{35}{37}{Translation.raise_edgeCount}

\lean{Copying/Translation}{39}{43}{Translation.pad_edgeCount}

\lean{Copying/Translation}{45}{47}{Translation.castDepth_edgeCount}

\lean{Copying/Translation}{49}{51}{Translation.padTo_edgeCount}

The required allowance is zero for variables, atoms, and constants;
one more than the body's allowance for abstraction; unchanged for unary
operations; and the maximum of the two subterm allowances for application
and binary operations. Induction on terms shows that variable and atom
renaming preserves this number, including renaming under a binder. This
is the depth part of a renaming law only; graph-level compatibility with
general environment renaming remains a separate obligation.

\lean{Copying/Translation}{54}{58}{Translation.depth}

\lean{Copying/Translation}{61}{65}{Translation.depth_rename}

\subsection{The full term translation}
\termdef{complete-term-translation}

Translate recursively by the seven scoped term constructors. The leaf
cases use the literal variable, atom, and constant clauses. Abstraction
boxes the recursively translated body; unary operations use the unary
clause. Application and binary operations first pad both recursive graphs
to the maximum of their allowances, then use their Figure~16 constructors.
The result type states the exact interface
$\star\to\star^k\otimes\diamond^h$ for every term in scope $(k,h)$.

Induction on the same syntax proves that every result is a language graph.
Each leaf uses its clause theorem, abstraction and unary operations use
their preservation laws, and the two-argument cases also use padding
invariance before their preservation laws. The generic implications from
the language predicate give global focus freedom and hole freedom. Thus
these properties hold recursively inside every translated box, with no
well-formedness assumptions on particular translated subterms left open.

\lean{Copying/Translation}{68}{79}{Translation.term}

\lean{Copying/Translation}{82}{92}{Translation.term_language}

\lean{Copying/Translation}{94}{94}{Translation.term_focusFree}

\lean{Copying/Translation}{96}{96}{Translation.term_holeFree}

The named-environment entry point accepts the original ordered distinct
lists and a term scoped by their lengths. It invokes exactly the same
translation; it does not restrict the environments to be empty or identify
variables with atoms. The two proposition graph definitions apply this
translation directly to the existing, unchanged \texttt{Term.privateCell}
and \texttt{Term.constantResult}. Their global language properties are
instances of the general theorem. Neither graph definition contains an
assumed observational equivalence or an alternate program.

\lean{Copying/Translation}{99}{100}{Translation.inEnvironment}

\lean{Copying/Translation}{103}{103}{Translation.privateCell}

\lean{Copying/Translation}{106}{106}{Translation.constantResult}

\lean{Copying/Translation}{108}{108}{Translation.privateCell_language}

\lean{Copying/Translation}{110}{110}{Translation.constantResult_language}

\subsection{Permutation, sharing, and full-program regressions}

For two variables and one atom in each subterm, the permutation sends
positions $(0,1,2,3,4,5)$ to $(0,1,4,2,3,5)$. This checks the required
movement of the left atom past the right variables. An application of two
constant graphs in this environment has eighteen outer edges: one operation,
four in each constant graph, and nine in the three two-use canonical trees.
The same count is checked for every binary operation, including both store
operations. These kernel-checked calculations would detect omission of the
two-use distributor edges.

\lean{Copying/FullTranslationExamples}{7}{11}{FullTranslationExamples.binary_environment_permutation}

\lean{Copying/FullTranslationExamples}{14}{16}{FullTranslationExamples.application_environment_edges}

\lean{Copying/FullTranslationExamples}{19}{22}{FullTranslationExamples.binary_environment_edges}

The unchanged left program requires three nested abstraction levels,
including the paper's let and sequencing sugar; the right requires two.
Their closed translated graphs have respectively five and two outer edges.
On the left these are the application, lambda, box, reference-creation,
and numeral edges. On the right they are lambda and box. The remaining
body graphs are inside boxes, not erased. For every named distinct-list
environment, a single theorem gives focus freedom and hole freedom of both
actual translated programs. It asserts structural well-formedness only;
no statement of termination equivalence follows from it alone.

\lean{Copying/FullTranslationExamples}{25}{25}{FullTranslationExamples.privateCell_depth}

\lean{Copying/FullTranslationExamples}{27}{27}{FullTranslationExamples.constantResult_depth}

\lean{Copying/FullTranslationExamples}{30}{31}{FullTranslationExamples.privateCell_outer_edges}

\lean{Copying/FullTranslationExamples}{34}{35}{FullTranslationExamples.constantResult_outer_edges}

\lean{Copying/FullTranslationExamples}{38}{44}{FullTranslationExamples.proposition12_2_graphs_wellFormed}

A self-application term provides a structural regression for syntax that
can express divergence. The generic language theorem checks its translation;
no evaluation or divergence theorem is claimed. A separate application
of an identity abstraction to unit has subterms of unequal box depth.
Its translated outer graph has four edges, checking that common-depth
padding has added no outer structure. Operational regressions for divergent
callbacks and repeated closure use still require the unfinished machine.

\lean{Copying/FullTranslationExamples}{47}{48}{FullTranslationExamples.omegaTerm}

\lean{Copying/FullTranslationExamples}{50}{51}{FullTranslationExamples.omega_translation_language}

\lean{Copying/FullTranslationExamples}{54}{55}{FullTranslationExamples.unequal_depth_application_edges}

The full build and all three axiom audits pass with only standard Lean
axioms. Construction of all Figure~16 clauses and the recursive term
translation is complete. The proposition terms and their desugaring remain
unchanged. The remaining correspondence, context, machine, and robustness
obligations are listed in Appendix~\ref{app:remaining-work}.
